
%%%===========================================
% \phantomsection
% \addcontentsline{toc}{section}{Abstract}

% =============================
% \section*{Abstract}
% =============================

\textbf{Environmental degradation often results from the socially burdensome trade-off between boosting productivity and overexploiting natural resources.} Two primary modes of resource use dominate this dynamic: extractivism and predation. Extractivism involves the removal of natural resources in a manner that, although potentially impactful, can ostensibly occur within regenerative or sustainable limits. By contrast, predation refers to the indiscriminate and often illegal depletion of renewable or non-renewable resources, such as through the burning of public lands or the extraction of minerals on indigenous territories, typically disregarding environmental laws and property rights.

These tensions between economic agents and environmental advocates permeate broader legal, ethical, and ecological paradigms that are difficult to stabilize socially and economically. As a result, the use of shared natural resources often evolves into a condition of imperfect competition, where productive agents, environmental organizations, and policymakers each act to optimize their interests at the expense of collective welfare.

\textbf{This research introduces an alternative perspective to reconciling environmental preservation and economic interests through the concept of \emph{economic commensalism}, which frames economic relations around rational neutrality in the use of natural resources. By repositioning private incentives to align with collective environmental goals, the framework fosters voluntary adherence to sustainable practices and aims to stabilize long-term ecological and economic outcomes through rational, market-based mechanisms.} Unlike traditional models that address environmental externalities primarily through corrective tools such as Pigouvian taxes, Coasian bargaining, Pareto efficiency, and others, or the prevailing institutional applied, based wether coercion wether sparse, inefficient, and unstructured positive incentives, this approach seeks to internalize ecological sustainability within the economic decision-making process itself.

\textbf{This thesis is structured around three complementary components:} \textbf{(i) a theoretical framework}, which employs intertemporal optimization methods, uncertainty management techniques, and dynamic learning processes to model decision-making behavior among economic agents; \textbf{(ii)} \textbf{mathematical modeling}, which formalizes the proposed mechanisms; \textbf{and (iii) empirical validation}, which tests their practical feasibility by applying these models to Brazil’s Mato Grosso region, an area marked by deforestation and agricultural expansion. \textbf{The research thus seeks to contribute to both economic theory and policy design by proposing a market-driven mechanism whereby economic actors voluntarily engage in sustainable practices while preserving economic returns and long-term viability.}

\newpage

% fim %